\section{End of Project Platform Evaluation}
\label{Evaluation:PlatformEvaluation}
\subsection{Initial Informal thoughts}
To begin my evaluation I first responded to some initial thoughts I received from peers shortly after they first used my project.

\begin{enumerate}
    \item " This feels like Obsidian meets Google Docs, but more student-friendly. Love how I can jump between whiteboards and text notes seamlessly!"
    \item "I didn't realize how much I needed the linking feature until I started using it. It's like my brain, but visualized. Can't wait to use this for revision season."
    \item "Honestly the search and tagging system is amazing. I don't Everything just feels connected."
    \item "Was playing around with the whiteboard mode - feels great. It's nice being able to sketch diagrams right in the app without switching tools."
    \item "This could legit replace like three different tools I'm using right now. Especially love how I can drop in resources from YouTube and other sites."
    \item "The visual graph of notes is honestly the most helpful thing I've used all year. Makes me feel like my revision has structure now."
\end{enumerate}

This initial feedback was overwhelmingly positive, with users appreciating the seamless integration of various features and the intuitive design. The linking and tagging systems were particularly highlighted as transformative for their study habits. I am so proud to see ideas I have implemented from the beginnings of this project being deemed highly logical and crucial to my peers.

\subsection{Conversations with users}
To dive further into the user experience, I conducted informal conversations with several users. The discussions focused on their overall impressions, specific features they found useful, and any areas for improvement. Please find below a list of notes I made from each interview:

\subsubsection*{User 1:}
Use 1 is a peer / friend of mine attending university currently. They were on a video call with me as they tried the software
The user expressed that the platform felt intuitive and easy to navigate. They particularly liked the ability to switch between text and whiteboard modes without losing context. 

Their main use case for the note-taking software is to record thoughts and ideas during university lectures which they feel this works well for however, they suggested that having an integrated system to transcribe an audio file into a note (e.g. from a lecturer speaking) and summarise / nicely format in my styling could be an amazing feature for further optimising user efficiency.

This is definitely a feature I can look to implement, after the EPQ program has concluded I definitely aim to continue development of the platform and will add this to my research list (\nameref{Planning:PostEPQPlan}).

He has continued to use my platform since our call and has now recorded over 50 notes!
\subsubsection*{User 2:}
User 2 is a business consultant who used the app for the first time whilst sitting down with me. They found the interactive tutorial very useful to quickly learn the platform and appreciated the help hints placed throughout the UI. They did however express a want for greater text/written instruction manuals that could they could read to further understand the more complex tools such as the maths blocks within notes.

They thought that the project would be very useful for their line of work as they could "weave together their notes from different projects in a succinic way"

\subsubsection*{User 3:}
User 3 is an Executive officer of a financial services organization who had a self-guided experience of the app. They appreciated the mathematics mode for its simple access to the vast array of tools the type-setting interface provides. T were glad I didn't "reinvent the wheel" and that the system was resultingly "simple to use, much more so then word which is inefficient as it is too reliant on clicking around". They also liked the whiteboard mode for its ability to freely record information in a huge (infinite) space


for more evaluation of my project please visit the End of Project Evaluation within my PPR Chapter beginning this document and additionally the presentation evaluation within the presentation chapter 


